\documentclass[10pt,a4paper]{amsart}
\usepackage[margin=19mm]{geometry}
\usepackage{amsmath,amssymb,mathtools}
\usepackage[scr=rsfs,cal=euler]{mathalfa}
\usepackage{xfrac}
\usepackage[unicode,hidelinks,bookmarks=false]{hyperref}
\newcommand{\ie}{i.e.\ }
\newcommand{\cf}{cf.\ }
\newcommand{\resp}{respectively\ }
\newcommand{\Subsection}{\subsection}
\providecommand{\nobreakdash}{\nobreak\hbox{-}}
\setlength{\emergencystretch}{2em}
\multlinegap0pt
\usepackage{amssymb}
\usepackage{mathtools}
\newtheorem{proposition}{Proposition}
\title{On spherical Milnor Classifying Spaces I:\\ differential geometry}
\author{Jean-Pierre Magnot}
\date{Source: arXiv:2605.16369v1 (CC BY 4.0)}
\begin{document}
\maketitle

\noindent\emph{Source and licence.} On spherical Milnor Classifying Spaces I:\\ differential geometry. Version arXiv:2605.16369v1 (CC BY 4.0), distributed by arXiv under the Creative Commons Attribution 4.0 International licence, http://creativecommons.org/licenses/by/4.0/, as recorded in the arXiv licence metadata for this exact version and shown on the abstract page https://arxiv.org/abs/2605.16369v1. Full licence notice: \texttt{licenses/arxiv-2605.16369-CC-BY-4.0.txt}.\par\smallskip

\noindent\textit{Editorial scope.} Counted unit: the article's proposition proposition (source0.tex lines 2192--2219). The article's complete original proof of it is delivered with it (source0.tex lines 2221--2324, one \texttt{proof} environment of 104 lines). No mathematics is written, repaired or re-derived: the statement and the proof are the article's own lines, in the article's own notation, and no retained line is substituted or re-flowed.  No further source-local context is needed: every cross-reference of every retained line resolves inside the two delivered spans.  Nothing was omitted from any retained span, so the package carries no omission record.  No retained line contains a citation, so the wrapper carries no bibliography and no bibliography entry.  No retained line contains a figure environment, an image-inclusion command or drawing code, so no image is delivered and the package declares no asset dependency.  Editorial formatting only, in addition: an AMS \texttt{amsart} preamble that restates the article's own declarations and macros used by the retained lines, namely the environments \texttt{proposition} on separate counters, together with the standard packages those lines need.  The retained environments print in this document as Proposition 1 in order of appearance, whereas the article prints the same items with the numbering of its own full document.  The complete native source of the article as distributed in the arXiv e-print for this exact version is bundled unchanged as \texttt{sources/200714/source0.tex}.
\begin{proposition}
Let \(\mathcal H\to E^{\mathrm{sph}}(G)\) be the effective horizontal
pseudo-bundle endowed with the non-degenerate metric \(\bar g\). Assume that:
\begin{enumerate}
\item on each finite local model \(S_I\times G^I\), the bundle
\(\mathcal H\) is represented by a finite-dimensional Riemannian vector bundle;
\item the transition maps preserve \(\bar g\);
\item \(\mathcal E\to E^{\mathrm{sph}}(G)\) is a Clifford module associated
with \(\mathrm{Cl}(\mathcal H)\);
\item the local Clifford modules glue compatibly with the transition maps.
\end{enumerate}
Then there exists a connection
\[
\nabla^{\mathcal E}:\Gamma(\mathcal E)
\longrightarrow
\Gamma(\mathcal H^*\otimes \mathcal E)
\]
compatible with the Clifford structure, i.e.
\[
\nabla^{\mathcal E}_X(c(v)s)
=
c(\nabla^{\mathcal H}_X v)s
+
c(v)\nabla^{\mathcal E}_X s
\]
for all local sections \(X,v\) of \(\mathcal H\) and all local sections
\(s\) of \(\mathcal E\).
\end{proposition}

\begin{proof}
We proceed locally and then check compatibility under gluing.

Let \(S_I\times G^I\) be one of the finite-dimensional local models of
\(E^{\mathrm{sph}}(G)\). On this model, the horizontal pseudo-bundle
\(\mathcal H\) is represented by a finite-dimensional vector bundle endowed
with the non-degenerate metric \(\bar g\). Therefore, by the ordinary
fundamental theorem of Riemannian geometry, there exists a unique
Levi-Civita connection
\[
\nabla^{\mathcal H,I}
\]
on \(\mathcal H|_{S_I\times G^I}\), characterized by
\[
\nabla^{\mathcal H,I}\bar g=0
\]
and by the vanishing of torsion.

This connection induces canonically a connection on the Clifford algebra
bundle
\[
\mathrm{Cl}(\mathcal H)|_{S_I\times G^I}.
\]
Indeed, for a local section \(v\) of \(\mathcal H\), one defines
\[
\nabla^{\mathrm{Cl},I}_X c(v)
=
c(\nabla^{\mathcal H,I}_X v),
\]
and extends this rule to products by the Leibniz rule. This is well-defined
because \(\nabla^{\mathcal H,I}\) preserves the metric. To see this, recall
that the Clifford relation is
\[
c(v)c(w)+c(w)c(v)=-2\bar g(v,w).
\]
Applying \(\nabla^{\mathrm{Cl},I}_X\) gives
\[
c(\nabla_X v)c(w)+c(v)c(\nabla_X w)
+
c(\nabla_X w)c(v)+c(w)c(\nabla_X v)
=
-2X(\bar g(v,w)).
\]
Since \(\nabla^{\mathcal H,I}\bar g=0\), the right-hand side equals
\[
-2\bar g(\nabla_X v,w)-2\bar g(v,\nabla_X w),
\]
which is exactly the derivative of the Clifford relation. Hence the
connection preserves the Clifford algebra structure.

Now let \(\mathcal E|_{S_I\times G^I}\) be the corresponding local
Clifford module. By assumption, \(\mathcal E\) is associated to the
Clifford structure. Locally, one may choose a connection
\[
\nabla^{\mathcal E,I}
\]
on \(\mathcal E|_{S_I\times G^I}\) compatible with
\(\nabla^{\mathrm{Cl},I}\), meaning that
\[
\nabla^{\mathcal E,I}_X(c(a)s)
=
c(\nabla^{\mathrm{Cl},I}_X a)s
+
c(a)\nabla^{\mathcal E,I}_X s
\]
for every local Clifford algebra section \(a\) and every local module
section \(s\). In particular, for \(a=v\in\Gamma(\mathcal H)\), this gives
\[
\nabla^{\mathcal E,I}_X(c(v)s)
=
c(\nabla^{\mathcal H,I}_X v)s
+
c(v)\nabla^{\mathcal E,I}_X s.
\]

It remains to check that these local connections glue. Let
\(S_J\times G^J\) be another local model and suppose the two charts overlap
through a refinement or inclusion of supports. By hypothesis, the transition
maps preserve the metric \(\bar g\), hence they identify the corresponding
Levi-Civita connections:
\[
\nabla^{\mathcal H,I}
=
\nabla^{\mathcal H,J}
\]
on the overlap, after transport by the transition map. Consequently the
induced Clifford algebra connections also agree on overlaps.

Since the Clifford modules are assumed to glue compatibly with the same
transition maps, the local module connections
\(\nabla^{\mathcal E,I}\) and \(\nabla^{\mathcal E,J}\) agree on overlaps.
Therefore they define a global connection
\[
\nabla^{\mathcal E}:\Gamma(\mathcal E)
\to
\Gamma(\mathcal H^*\otimes\mathcal E).
\]

The compatibility identity is local, and it has been verified on every
finite-dimensional local model. Since the diffeology on
\(E^{\mathrm{sph}}(G)\) is generated by these local models, the identity
holds globally. Hence \(\nabla^{\mathcal E}\) is a Clifford-compatible
connection.
\end{proof}

\end{document}
